\chapter{Decomposition into components}\label{ch:decomp}
\chapterbrief
{the operator is known and factored: $P=c\prod_{j=1}^sQ_j$ with $Q_j=P_j^{m_j}$ pairwise coprime. For the main theorem $\Omega$ is open, connected and simply connected (more precisely $H^1_{dR}(\Omega)=0$).}
{once $P$ is known, we want to split $u$ into the parts generated by each piece of $P$, so that the pieces can be studied separately.}
{second link of the chain. The split exists for \emph{factors} of the symbol, with a finite and exactly known ambiguity. It fails for additive \emph{terms}, and that failure is what Chapter~\ref{ch:interaction} measures.}

\section{Which pieces? Terms versus factors}

An operator can be broken up in two different ways. Take the heat operator
\[
L=\partial_t-D\partial_x^2 .
\]
Its \emph{terms} are $L_1=\partial_t$ and $L_2=-D\partial_x^2$: they are added. Its symbol $\tau-D\xi^2$ is irreducible, so it has only one \emph{factor}. For the wave operator the situation is the opposite: $\tau^2-\xi^2$ has two terms and also two factors, $\tau-\xi$ and $\tau+\xi$. The first question is which kind of piece gives a decomposition of the solution.

\subsection{Terms do not give components}

\begin{proposition}{Splitting by terms fails \led{R2-T3} \proved}\label{prop:terms}
Let $L=L_1+L_2$ with constant coefficients. If $u=a+b$ with $L_1a=0$ and $L_2b=0$ solves $Lu=0$, then $u$ is also killed by $L_1L_2$:
\[
\ker L\cap(\ker L_1+\ker L_2)\ \subseteq\ \ker(L,\ L_1L_2).
\]
The common zero set of $L$ and $L_1L_2$ is $\Var(L_1)\cap\Var(L_2)$. When this set is finite, only a finite-dimensional space of exponential-polynomial solutions can be split by terms.

\dimref{R2-T3}
\end{proposition}

\begin{intuition}
The two operators commute, so
\[
L_1L_2\,u=L_2(L_1a)+L_1(L_2b)=0 .
\]
A solution split by terms satisfies one extra equation. A generic solution does not.
\end{intuition}

\begin{example}{The heat equation cannot be split by terms}
Take $L_1=\partial_t$ and $L_2=-D\partial_x^2$. The kernels are
\[
\ker L_1=\{f(x)\},\qquad \ker L_2=\{a(t)+b(t)\,x\}.
\]
Every element $f(x)+a(t)+b(t)x$ of the sum is killed by $\partial_t\partial_x^2$. A heat solution is not. For $u=e^{-D\kappa^2t}\cos\kappa x$:
\[
\partial_x^2u=-\kappa^2u,\qquad \partial_t\partial_x^2u=D\kappa^4u\neq0 .
\]
So a generic heat solution is \emph{not} a sum of a solution of $u_t=0$ and a solution of $u_{xx}=0$.

The exceptions are few. The function $u=\tfrac12x^2+Dt$ solves the heat equation and lies in $\ker L_1+\ker L_2$, but it is a polynomial: the exceptional space is finite-dimensional.
\end{example}

\subsection{Factors do give components}

\begin{example}{The wave equation and d'Alembert}\label{ex:wave}
Let $c>0$ and
\[
P=\tau^2-c^2\xi^2=(\tau-c\xi)(\tau+c\xi)=Q_1Q_2 .
\]
A direct check gives
\[
(\partial_t-c\partial_x)\,f(x+ct)=0,\qquad (\partial_t+c\partial_x)\,g(x-ct)=0 .
\]
So $\ker Q_1=\{f(x+ct)\}$ (profiles moving left) and $\ker Q_2=\{g(x-ct)\}$ (profiles moving right). D'Alembert's formula says that every solution is a sum of the two:
\[
\ker P=\ker Q_1+\ker Q_2 .
\]
The pieces are unique up to a constant: $f\mapsto f+k$, $g\mapsto g-k$. This ambiguity has dimension $1$, and it equals $\dimc\ring/(Q_1,Q_2)$:
\[
Q_1+Q_2=2\tau,\qquad Q_2-Q_1=2c\xi\quad\Longrightarrow\quad(Q_1,Q_2)=(\tau,\xi),\qquad \ring/(\tau,\xi)\cong\C .
\]
\end{example}

\begin{figure}[!htb]
\centering
\begin{tikzpicture}
\begin{axis}[width=7.4cm,height=4.9cm,xlabel={$x$},ylabel={$u$},xmin=-6,xmax=6,ymin=-0.05,ymax=1.15,legend style={font=\scriptsize,at={(0.98,0.97)},anchor=north east,draw=none,fill=none},
 title={\small $u=f(x+t)+g(x-t)$},domain=-6:6,samples=100,no markers]
\addplot[cG,thick] {1/(1+(x+2)^2)+exp(-(x-2)^2)};
\addplot[cA,thick,dashed] {1/(1+(x+0)^2)+exp(-(x-0)^2)};
\addplot[cB,thick] {1/(1+(x-2)^2)+exp(-(x+2)^2)};
\legend{$t=-2$,$t=0$,$t=2$}
\end{axis}
\end{tikzpicture}
\hfill
\begin{tikzpicture}[scale=0.9]
\draw[->] (-2.6,0)--(2.8,0) node[right]{$x$}; \draw[->] (0,-0.1)--(0,3) node[above]{$t$};
\foreach \a in {-2,-1.2,-0.4,0.4,1.2,2}{
  \draw[cA,thick] (\a,0)--(\a+2.2,2.6); \draw[cB,thick] (\a,0)--(\a-2.2,2.6);}
\node[font=\scriptsize\sffamily,text=cG] at (0,-0.6) {characteristics of $Q_1$ and $Q_2$};
\end{tikzpicture}
\caption{Left: the d'Alembert solution with $f(s)=1/(1+s^2)$ moving left and $g(s)=e^{-s^2}$ moving right ($c=1$). Right: the two families of characteristic lines, one per factor. Each factor of the symbol carries one rigidly translating component.}\label{fig:wave}
\end{figure}

\section{Why factors work: the Chinese remainder theorem}

\textbf{One variable.} If two polynomials $Q_1,Q_2$ in one variable are coprime, there are polynomials $a,b$ with
\[
a\,Q_1+b\,Q_2=1 .
\]
Then every $u$ with $Q_1Q_2u=0$ splits explicitly as
\[
u=\underbrace{b\,Q_2\,u}_{\in\ker Q_1}+\underbrace{a\,Q_1\,u}_{\in\ker Q_2},
\]
and the sum of the two kernels is direct.

\textbf{Two variables.} ``Coprime'' now means only that the common zeros are finitely many (B\'ezout), not that $(Q_1,Q_2)=\ring$. The finite quotient $\ring/(Q_1,Q_2)$ measures the failure of the one-variable picture.

\begin{lemma}{Finite quotient \led{C12-L3} \proved}\label{lem:C12L3}
Let $f,g,h\in\ring\setminus\{0\}$.
\begin{enumerate}[label=(\alph*)]
\item If $f,g$ are coprime, the ideal $(f,g)$ contains a nonzero polynomial in $\xi$ alone and a nonzero polynomial in $\tau$ alone. Hence $\dimc\ring/(f,g)<\infty$.
\item If $f$ is coprime to $g$ and to $h$, then
\[
\dimc\ring/(f,gh)=\dimc\ring/(f,g)+\dimc\ring/(f,h).
\]
\item $(f,g)=\ring$ if and only if $\Var(f)\cap\Var(g)=\emptyset$.
\end{enumerate}
\dimref{C12-L3}
\end{lemma}

The geometric meaning of $\dimc\ring/(Q_i,Q_j)$ is the number of crossing points of the curves $\Var(Q_i)$ and $\Var(Q_j)$, counted with intersection multiplicity.

\section{The decomposition theorem}

\begin{theorem}{Factor decomposition \led{C12-T1, C12-T1$'$} \proved}\label{thm:factor}
Let $\Omega\subset\R^2$ be open, connected, with $H^1_{dR}(\Omega)=0$ (for example simply connected: a rectangle, a convex set, a slit disc). Let $\mathcal A$ be $\mathcal D'(\Omega)$ or $C^\infty(\Omega)$, and $P=c\prod_{j=1}^sQ_j$ with pairwise coprime $Q_j$. Then:
\begin{enumerate}[label=(\roman*)]
\item every solution splits:
\[
\ker_{\mathcal A}P=\ker_{\mathcal A}Q_1+\dots+\ker_{\mathcal A}Q_s ;
\]
\item the ambiguity (the $s$-tuples of components with sum $0$) has dimension exactly
\[
\delta(P)=\sum_{i<j}\dimc\ring/(Q_i,Q_j),
\]
and consists of exponential-polynomials;
\item the splitting is unique if and only if $\Var(Q_i)\cap\Var(Q_j)=\emptyset$ for all $i\ne j$.
\end{enumerate}
\dimref{C12-T1p}
\end{theorem}

\begin{intuition}
At a crossing point $z$ of two factor curves, the exponential $e^{\langle z,(t,x)\rangle}$ belongs to both kernels. It can be moved freely from one component to the other. The number of such free moves, counted with multiplicity, is $\delta(P)$.

Algebraically the proof uses the exact sequence
\[
0\to\ring/(Q_1Q_2)\to\ring/(Q_1)\oplus\ring/(Q_2)\to\ring/(Q_1,Q_2)\to0
\]
and shows that applying $\mathrm{Hom}(\,\cdot\,,\mathcal A)$ keeps it exact. The only possible obstruction lives on the finite-dimensional module $\ring/(Q_1,Q_2)$, and it vanishes when $\Omega$ has no holes (Lemma~\ref{lem:twisted}).
\end{intuition}

\begin{example}{The ambiguity on four operators}
\begin{center}\small
\begin{tabular}{@{}L{3.0cm}L{4.4cm}L{1.0cm}L{4.4cm}@{}}
\toprule
Operator & Factors & $\delta(P)$ & Ambiguity\\\midrule
wave $\tau^2-\xi^2$ & $\tau-\xi,\ \tau+\xi$ & $1$ & constants\\
beam $\tau^2+\xi^4$ & $\tau-i\xi^2,\ \tau+i\xi^2$ & $2$ & $\spn\{1,x\}$\\
heat--transport & $\tau-\xi^2,\ \tau-\xi$ & $2$ & $\spn\{1,\ e^{t+x}\}$\\
parallel transports & $\tau-\xi,\ \tau-\xi-1$ & $0$ & none (unique)\\\bottomrule
\end{tabular}
\end{center}
For heat--transport the curves $\tau=\xi^2$ and $\tau=\xi$ meet at $(0,0)$ and $(1,1)$, so the shared exponentials are $1$ and $e^{t+x}$. For the two parallel transports the lines $\tau=\xi$ and $\tau=\xi+1$ never meet, and the splitting is unique.
\end{example}

The key analytic tool is a version of the Poincar\'e lemma ``closed forms are exact'' with an exponential weight.

\begin{lemma}{Twisted Poincar\'e lemma \led{C12-L1$'$, L2$'$} \proved}\label{lem:twisted}
Let $\Omega$ be connected with $H^1_{dR}(\Omega)=0$, and $z=(z_1,z_2)\in\C^2$. Write $D_t=\partial_t-z_1$ and $D_x=\partial_x-z_2$. If $\varphi,\psi\in\mathcal A$ satisfy the compatibility condition
\[
D_x\varphi=D_t\psi,
\]
then there is $w\in\mathcal A$ with $D_tw=\varphi$ and $D_xw=\psi$. Consequently, for every finite-dimensional $\ring$-module $M$, the obstruction group $\mathrm{Ext}^1_\ring(M,\mathcal A)$ vanishes.

\dimref{C12-L1p}
\end{lemma}

Multiplying by $e^{-z_1t-z_2x}$ turns the twisted derivatives into ordinary ones, and the lemma becomes the classical statement that a closed $1$-form on a simply connected region is exact.

\section{Where the decomposition fails: holes}

\begin{proposition}{Holes break it \led{C12-P10} \proved}\label{prop:holes}
Let $\Omega$ contain a closed curve $\gamma$ that winds around a point $\alpha\notin\Omega$. For the Laplacian
\[
P=\tau^2+\xi^2=(\tau-i\xi)(\tau+i\xi),
\]
the harmonic function $u=\log|(t,x)-\alpha|^2$ is \emph{not} a sum of a solution of $\partial_t-i\partial_x$ and a solution of $\partial_t+i\partial_x$.

\dimref{C12-P10}
\end{proposition}

\begin{figure}[!htb]
\centering
\begin{tikzpicture}
\fill[cA!10] (0,0) circle (1.7); \fill[white] (0,0) circle (0.5);
\draw[cA,thick] (0,0) circle (1.7); \draw[cA,thick] (0,0) circle (0.5);
\node at (0,0) {\small $\alpha$};
\draw[cB,thick,-{Stealth}] (1.1,0) arc (0:340:1.1);
\node[cB,font=\scriptsize] at (0,-2.05) {$\oint_\gamma d\arg z=2\pi$};
\node[align=left,font=\small,anchor=west] at (2.4,0) {$u=\log|z-\alpha|^2$ is harmonic on the annulus,\\ but $u\neq h(z)+\overline{g(z)}$ with $h,g$ single-valued\\ holomorphic: the logarithm has monodromy.};
\end{tikzpicture}
\caption{The Laplacian factors into two Cauchy--Riemann operators, but on an annulus the factor decomposition fails. The obstruction is topological: $H^1_{dR}(\Omega)\neq0$.}\label{fig:hole}
\end{figure}

For a fixed $P$ the exact condition is cohomological.

\begin{theorem}{Cohomological criterion \led{OPF-T4} \pw}\label{thm:cohom}
Let $\Omega$ be connected, with first Betti number $b_1(\Omega)$. Let $\mathcal N$ be the space of $s$-tuples of global exponential-polynomial solutions with sum zero (dimension $\delta(P)$). The cokernel of the sum map
\[
\Sigma_\Omega:\ \bigoplus_j\ker_{\mathcal A(\Omega)}Q_j\to\ker_{\mathcal A(\Omega)}P
\]
is the kernel of an explicit map $H^1(\Omega;\C)\otimes\mathcal N\to\bigoplus_jH^1(\Omega;\mathcal K_j)$. So the decomposition holds on $\Omega$ iff that map is injective, and
\[
\dimc\operatorname{coker}\Sigma_\Omega\le\delta(P)\;b_1(\Omega).
\]
\dimref{OPF-T4}
\end{theorem}

\begin{example}{Unique splittings survive holes}
If $\delta(P)=0$ (disjoint factor curves), the bound gives $\operatorname{coker}\Sigma_\Omega=0$ on \emph{every} connected region, with or without holes. The two parallel transports $\tau-\xi$ and $\tau-\xi-1$ split on an annulus as well.
\end{example}

\section{Explicit projectors}

For evolution factors the components are given by a formula. Let
\[
P=\prod_{j=1}^N\big(\tau-a_j(\xi)\big)
\]
with distinct polynomials $a_j$. A component is an eigenfunction of $\partial_t$:
\[
\partial_tu_j=a_j(\partial_x)\,u_j .
\]
Apply $\prod_{k\ne j}(\partial_t-a_k(\partial_x))$ to $u=\sum_lu_l$. Every $u_l$ with $l\ne j$ is killed, because the product contains its own factor. The component $u_j$ is multiplied by $\prod_{k\ne j}(a_j-a_k)(\partial_x)$. Dividing gives the Lagrange formula
\begin{equation}\label{eq:lagrange}
u_j=\prod_{k\ne j}\frac{\partial_t-a_k(\partial_x)}{a_j(\partial_x)-a_k(\partial_x)}\;u .
\end{equation}
In Fourier variables the right-hand side is a multiplier that is rational in $\xi$, with poles where two $a_k$ coincide. This is Eq.~\eqref{eq:spectral} with the companion matrix replaced by $\partial_t$.

\begin{example}{Two transports}
Let $a_1(\xi)=\xi$ and $a_2(\xi)=\xi+1$. The formula gives
\[
u_1=\frac{\partial_t-\partial_x-1}{-1}\,u,\qquad u_2=\frac{\partial_t-\partial_x}{1}\,u .
\]
These are differential operators, with no poles, because $a_1-a_2=-1$ never vanishes: the two lines are disjoint and the splitting is unique.
\end{example}

\begin{proposition}{Projectors and their poles \led{II-T3, II-T4} \proved}\label{prop:poly-proj}
The components are obtained by \emph{polynomial} operators if and only if the factor curves are pairwise disjoint, $(Q_i,Q_j)=\ring$. Only factors, not terms, admit projectors that act on one value of $u$ at a time.

\dimref{II-T3}
\end{proposition}

\begin{keypoint}
A ``component'' of $u$ belongs to a \emph{factor} of the symbol, not to a term of the operator. On a region without holes the factor split always exists, and the overlap between components is finite-dimensional and measured exactly by $\delta(P)$. The split by terms fails except on a finite-dimensional space.
\end{keypoint}
